\section{A common normalization of the point readings}

Fix a rank index $i$ and write $V=V_c(D,i)$, $W=W_c(D,i)$,
$n=\operatorname{rank} V$ and $R$ for the fixed real datum. Use the
original core, its weight and vector ties, and the extension of its
joint-harmonic map (\texttt{IoExt}). The tuple set, the coefficient
field of a tuple record $O$, and its class representatives may vary.
The coefficient map on the original lower field is fixed. The issue
is to choose a function of the automorphic point before any of these
records have been specified.

% BEGIN P74I1
\paragraph{P74I1. Comparing two projectively equivariant readings.}
Let $M$ be a nonzero irreducible finite-dimensional complex
representation of a group $\Gamma$, generated by elements with
$\operatorname{tr}\rho(\gamma)\ne0$. Let $A,B:M\to N$ be linear.
Suppose $A\rho(\gamma)=S_\gamma A$ and
$B\rho(\gamma)=c_\gamma S_\gamma B$, with linear $S_\gamma$
on $N$. Suppose also that $Aw=Bz\ne0$. Both kernels are invariant
and therefore zero. Since $\rho(\gamma)$ is invertible,
$c_\gamma\ne0$. The subspace $A^{-1}(B(M))$ contains $w$ and
is invariant since $S_\gamma B(y)=B(c_\gamma^{-1}\rho(\gamma)y)$.
It follows that $T=B^{-1}A$ is defined on all of $M$ and is invertible.
Conjugation by $T$ carries $\rho(\gamma)$ to
$c_\gamma^{-1}\rho(\gamma)$. Taking traces at the stated
generators gives $c_\gamma=1$ there. Commutation with $T$ extends
to their products and inverses. Any nonzero eigenspace of $T$ is now
invariant, so $T=aI$. In particular $A=aB$ and $aw=z$.

Here $\Gamma=U(V)(F)$ acts through the compact archimedean group of
the definite space $V$. Its dense image can be seen in the Cayley
chart at the identity. The skew-Hermitian Lie algebra is an $F$-vector
space and weak approximation gives simultaneous rational
approximation at the archimedean places. The chart
$X\mapsto(1-X)(1+X)^{-1}$ therefore puts a neighborhood of the
identity in the closure of the rational group. The product of the
compact unitary groups is connected, hence the closure is the entire
group. The finite action on this coefficient module is trivial. Thus
an invariant rational subspace over $\mathbb C$ is invariant under
the real group by continuity and is zero or all of $M$ by its given
irreducibility. The elements of nonzero trace include the preimage
of an open identity neighborhood, since the trace at $1$ is
$\dim M>0$. For clarity, take the subgroup generated by this
preimage. Its image closure contains an open identity neighborhood
and hence the full connected group. Every rational element then
differs from an element of the subgroup by an element mapping into
the original neighborhood. This proves the required abstract
generation statement as well as density.
% END P74I1

% BEGIN P74I2
\paragraph{P74I2. One unit comparison on a fixed original record.}
Write $v_t$ for the stored theta datum, and $P\mapsto E_t(P)$ for
application of the creation polynomial to $v_t$ in the identity
source basis. At the base parameter $t_*=\mu_0$ this map is injective.
Indeed, the creation and annihilation operators on the Fock vacuum
satisfy the canonical commutators. In dual coordinates, applying
an annihilator to $E_{t_*}(P)$ applies the corresponding partial
derivative to $P$, up to a fixed nonzero normalization. The real
basis is orthonormal, so these coordinates give all derivatives.
A nonzero term of highest total degree can be differentiated to a
nonzero constant with its own multi-index. The image of $1$ is
nonzero. For the base finite datum take the rectangular identity on its
supported half. This is possible since $\operatorname{rank}W\le n$.
Its lower block vanishes and its top determinant is $1$ at every
residue prime. Both finite factors are nonzero. The vacuum has
unit norm, and the half-lattice Schrödinger realization on this
nonzero vacuum is injective (its distinct translates are separated
by lattice characters). Thus a polynomial in the kernel would have
all these constant derivatives in the kernel, an impossibility.

Fix an allowed identification of letters $\ell$ and the lowest-type
vector $u$ specified by the tie. Denote by $P_O(v)$ the complex
coefficient extension and renaming of $O.\mathrm{io}(v)$. On the
complex coefficient module $M$ form the two readings
\[
 A(w)=\mathrm{ptP}_{t_*}(\mathrm{ioP}(u,w)),\qquad
 B(O.c'_0 v)=E_{t_*}(P_O(v)).
\]
The second formula extends complex linearly. In fact the coefficient
tie identifies $O.c'_0$ on a basis extended from the lower field
with the original complex basis. The basis remains independent and
spanning on extending coefficients. The same argument applied to
the original rational unitary action shows that $O.c'_0$ intertwines
it with the action on $M$. This can be checked on the lower basis.

The extension and extremal-vector pins give a unit vector $w_0$ and
$\mathrm{ioP}(u,w_0)=\Delta$, the specified product of extremal
minors. The source basis in $A$ is the stored $B_V$.
Each $B_V$ is unitary as $V$ has identity Gram matrix. Take the
real group element $g$ with those components, and let $\chi$ be the
unit scalar in the real oscillator basis-change formula. Compact
equivariance, applied also to $g^{-1}$, gives
\[
 A(w_*)=E_{t_*}(\Delta),\qquad
 w_*=\rho(g^{-1})(\chi w_0),\qquad \|w_*\|=1 .
\]
Set $z_*=O.c'_0(O.\mathrm{hw})$. Its norm is one: the coefficient
pairing of the tensor of the pinned extremal vectors with itself in
the identity Gram matrix is $1$. The polynomial pin for $O.\mathrm{hw}$
is the product of the geometric minor polynomials. In the positive
embedding its row labels start at zero. In the conjugate embedding
they start after the $p$ positive rows. The weight tie gives the
successive highest-weight differences as their exponents. Under each
permitted renaming this gives the same $\Delta$, so
$B(z_*)=E_{t_*}(\Delta)\ne0$. Alternatively, nonvanishing follows
from injectivity of $O.\mathrm{io}$, of the letter renaming, and of
$E_{t_*}$, since $z_*\ne0$.

For $\gamma\in U(V)(F)$ let $g_\infty$ be its real part. The finite
part acts trivially on $M$. Compact equivariance gives
$A\rho(\gamma)=\omega_{t_*}(g_\infty,1) A$. For $B$, apply the
rational substitution pin to $P_O(v)$ and then the real oscillator
substitution formula. The remaining scalar is
$c_O(\gamma_E)\chi(g_\infty)^{-1}$, the pinned rational polynomial
factor divided by the real oscillator factor. The two identities on
the spanning set $O.c'_0(v)$ extend to $M$. P74I1 consequently gives
$A=a_OB$ and $a_Ow_*=z_*$. Both vectors have norm one, so $|a_O|=1$.
The latter equation is independent of the letter identification, and
$w_*\ne0$ makes its scalar unique. The vector $u$ is also fixed,
since each of its components is a subtype with prescribed underlying
vector. Finally strip off $E_{t_*}$ using its injectivity and the
source-basis substitution formula. This gives the polynomial identity
before applying any tuple vacuum, so applying $E_t$ gives
\[
 \mathrm{ptP}_t(\mathrm{ioP}(u,O.c'_0 v))
       =a_O E_t(P_O(v))\qquad\text{for every }t,v.
\]
This is $\mathrm{IoAt}(O,v_0,a_O)$. If no letter identification is
allowed the assertion is vacuous and $a_O=1$ works.
% END P74I2

% BEGIN P74I3
\paragraph{P74I3. A single coefficient phase over all extensions.}
Fix the original lower map $b=c_{0,L}(i)$. If no pinned record
satisfies coefficient base change for this $b$, use $a=1$.
Otherwise choose one such record $O_*$, allowing an arbitrary tuple
set and coefficient extension, and take $a=a_{O_*}$ from P74I2.
Let $O$ be any other record satisfying both pins. Fix an allowed
letter map and lowest-type input. Apply its reading and the anchor
reading at the common base parameter and their respective extremal
vectors. Their left sides agree: the tensor $v_{\rm ext}$ of the
extremal vectors over the lower field maps to $O.\mathrm{hw}$ and
$O_*.\mathrm{hw}$ under the two field extensions, since coefficient
extension commutes with tensors and extremal minors. Both complex
coefficient vectors are consequently $b(v_{\rm ext})$.
On the right, the variables are renamed with the same map.
Coefficient extension of each determinant polynomial is the
determinant of the same matrix of indeterminates, and products and
powers commute with extension. Thus both right sides use precisely
$E_{t_*}(\Delta)$. It is nonzero by P74I2. Cancellation gives
$a_O=a$, and P74I2 then proves the full reading on $O$. This shows
$|a|=1$ and proves \texttt{exists\_scalar\_core} with its order of
quantifiers across all tuple sets. No coefficient field containing
both extensions needs to be chosen besides their specified
embeddings in $\mathbb C$.
% END P74I3

% BEGIN P74I4
\paragraph{P74I4. The cover value on every marked analytic fiber.}
Fix a pinned record $O$ and any scalar $c$ satisfying its original
$\mathrm{XiConst}(O,c)$. For an arbitrary tuple $t$, denote its
higher model by $U_2$. Let $k$ be the adelic element and $x_k$ its
stipulated determinant root in the test defining $\mathrm{CovAt}$.
We claim that at every framed point $y_2$ equipped with the prescribed
analytically marked fiber,
\[
              \mathrm{ev}(c_\delta,y_2)=c\,\xi_t(x_k).
\]
The higher dictionary specified in $\mathrm{XiConst}$ gives a
component, its frame, and a rational $\gamma$ such that the lattice
and frame there are those of $\gamma k$ and the marking is its
adelic level marking. The core fixes the CM type, orthonormal
reference frame and common Hermitian form. Write $D_\gamma$ for the
rational unitary matrix for $\gamma$. It maps the test lattice for
$k$ bijectively onto the component lattice for $\gamma k$ by the
finite-lattice formula, preserves the Hermitian form, and its
inverse carries the component frame back to the test frame. Period
labels agree after the permutation given by the two ordering maps.
The two level markings agree after this change: at each completed
lattice their difference lies in the level multiple, by the two
adelic marking congruences. A rational vector in this multiple at
every completion belongs to the level multiple of the rational
period lattice. Thus the resulting polarized complex tori agree as
marked moduli objects. Their points in the higher fine moduli space
are equal. The section $c_\delta$ has weight and index zero, so
its evaluation depends only on this common underlying point.

In the value law at that dictionary point a rational unitary factor
$\gamma'$ may occur, from the representatives in $\mathrm{XiConst}$.
Its determinant has a root $b\in E^\times$ by the rational
determinant-root law (Hilbert 90). The product of its diagonal
image with $x_k$ is a root for $\gamma'k$. The tuple pin reads the
dictionary character as $\xi_t(bx_k)$, which equals $\xi_t(x_k)$
by multiplicativity and triviality on $E^\times$. Substituting
in the given $\mathrm{XiConst}$ value gives the asserted formula
with the original $c$, including $c=0$. This proves
\texttt{covAt\_of\_xiConst\_core} for the whole marked fiber.
% END P74I4

% BEGIN P74I5
\paragraph{P74I5. The unit attached to each component.}
Take a pinned record at the base parameter. Its embedding on the
lower field is the distinguished one, and its lower dictionary,
character $\mu$ and frame are the corresponding ones of $T$. Fix a
component with period lattice $L$. Choose a symplectic basis of $L$,
compatible coordinate labels and a classical vacuum $v$ for
$\mathcal O_E^n\otimes L$. In these coordinates $\Theta(v)$ denotes
its classical theta function. There are two comparisons to make.

First let $\vartheta_x$ be the analytic reading of the geometric
generator in a unit frame $x$. It has the canonical multiplier and
Gaussian-weighted squared mean one. The $n$-fold product and
$\Theta(v)$ belong to the same vacuum line on the product torus.
Indeed, the Heisenberg lattice acts on both by the prescribed
semicharacter, and the holomorphic space for a principal polarization
is one-dimensional, also for its external product. Thus
$\prod_b\vartheta_x(z_b)=C(x)\Theta(v)(x,z)$. The vacuum mean
formula gives $|C(x)|=\rho_L>0$ independently of $x$: integrate
the squared norms with the same Gaussian metric, using the product
probability measure on the torus. Equivalently one can compute the
classical vacuum norm in a symplectic basis by orthogonality of its
Fourier modes followed by the Gaussian integral. The normalized
vacuum norm at a unit frame is independent of that frame.

Here $C(x)$ can be made constant on the component. To see the phase,
translate by the level division points and take degree-zero theta
values at the origin. These give pairs of holomorphic readings
$(S_q,T_q)$ of the same weight, with $S_q=C(x)T_q$ on unit frames.
At every unit frame some $T_q$ is nonzero. More explicitly the
specified half-level lattice $\mathfrak m L$ is self-dual for
the original alternating form divided by $M$, by
$\mathfrak m\overline{\mathfrak m}=(M)$ and self-duality of $L$.
On that lattice the original theta function is of degree $M\ge3$.
The evaluations on its $M$-division grid have trivial common
kernel. For a principal polarization the translations by
$M$-torsion act irreducibly on the degree-$M$ theta space:
in a symplectic basis the diagonal half separates the
$M^g$ residue Fourier modes, and the other half permutes them
transitively. Evaluation at zero is a nonzero functional since the
$M$th power of a principal theta bundle, with any flat twist, is
generated there. For example, take a product of $M$ translates of
the principal section, with prescribed sum of translation parameters
to realize the twist. For $M\ge2$ the parameters can be chosen
off the finitely many proper divisors imposed by vanishing of a
factor at zero. The common kernel of all translates of the evaluation
is invariant and proper, hence zero. The marked-root congruences
identify this grid with the required values modulo periods. The period multipliers are nonzero. Taking the same nonvanishing
coordinate in each column gives the product assertion.

Where $T_q\ne0$, the ratio is holomorphic and has constant absolute
value. The common weight passes this last assertion from unit frames
to arbitrary frames by triangular change of basis. It follows from
the open mapping theorem that the ratio is locally constant.
The local ratios agree on overlaps since they all give $C(x)$ on
unit frames and the common weight cancels. The period component is
connected, so there is a single $C\ne0$, with $|C|=\rho_L$.
Changing the basis and ordering only makes a simultaneous coordinate
change in the Hermitian form, the Gaussian, and the arguments of
$\vartheta$, so the comparison applies to every generator reading.

Second, if the component lattice occurs adelically, choose a finite
$k_0$ moving the base lattice to it. The vector
$\omega_{t_*}(1,k_0)\mathrm{vac}$ satisfies the moved lattice and
Fock equations. In the complex period coordinates these are exactly
the same vacuum-line equations for $v$: the finite character is
the integral semicharacter, and the Fock equation identifies the
real restriction with a Gaussian times an entire theta function.
Consequently the moved vector reads $\kappa_0 v$, with
$\kappa_0\ne0$ by the unit norm of the adelic vacuum. Its norm fixes
the modulus without any choice of a theta constant. Indeed let
$H_x$ be the frame Hermitian form and write
$F_x(z)=|\vartheta_x(z)|^2\exp(-\pi H_x(z,z))$. It is periodic
and its torus mean is $1$. In the real period coordinates the
squared norm of the adelic reading has mean
\[
   \frac{|\kappa_0|^2}{|C|^2}\,
      \operatorname{mean}_{z}\prod_{b=1}^{n}F_x(z_b)
       =\frac{|\kappa_0|^2}{|C|^2}.
\]
For the adelic inner product this change of domain is legitimate:
the finite part is invariant in squared norm under an open lattice.
Unfolding its finite cosets gives a real period fundamental domain,
with its probability Haar measure. The period-coordinate map takes
Haar probability to Haar probability, and Fubini evaluates the
product. Since the adelic norm squared is $1$, $|\kappa_0|=|C|$.

Set $\kappa=\kappa_0/\mu(k_0)$. A second finite mover $k_1$ to the
same lattice changes the moved vacuum by the stabilizer character
$\mu(k_1k_0^{-1})$. The real part of a general mover acts by
moving the frame and its determinant character. Thus at $h$ the
vacuum reads $\kappa\mu(h)v$. The base twist is trivial, so the
lower character is the one tied to the original Weil action.
Left multiplication by a rational unitary element leaves $\mu(h)$
unchanged. Therefore $C_k=\kappa/C$ is a unit for the component
$k$, and the moved adelic vacuum, on arguments of finite part zero,
is $C_k\mu(h)$ times the original Gaussian times the product of the
actual generator readings. On a lattice reached by no adelic point,
any unit serves. Choosing one unit per component proves
$\mathrm{CcIs}$ with all its frame, basis and reading quantifiers.
% END P74I5

% BEGIN P74I6
\paragraph{P74I6. The signed contribution of one split class.}
The isometry between the level lattice pairing and the trace of the
Hermitian pairing supplies a sign $\delta\in\{1,-1\}$ before any
tuple or representative is chosen. Fix now the test data of
$\mathrm{ClsAt}$: a level and residue function $\phi$ of the stated
type, a split-supported representative $h$, an arbitrary target
mover $k$, and roots $Y(q)$ of the given higher marking.
Write $\mathcal Q=(\mathrm{LevQ})^n$, $\widehat\phi$ for the
unnormalized finite Fourier transform, $S=\mathcal O_E^n$,
$\Lambda$ for the prescribed source class lattice, $r=\operatorname{rank}W$,
and $y$ for the split determinant parameter of $h$. Let $v$ be
the classical vacuum in the test and suppose the original target
move $\omega_t(1,k)\mathrm{vac}$ reads $b v$.
Set $I=[\Lambda:\Lambda\cap S]$ and $\chi_W$ equal to the second
character of the tuple's original theta pair.

On the original half-lattice the finite theta test is the function
$\phi$ of the reduced matrix of columns. The lattice embedding
identifies rational columns in the standard model with the original
level lattice. Any finite lattice input is congruent to such columns
modulo $M$ by density in its completion. The level hypothesis puts
$M$ in all the residue primes used by the datum, so both tests have
the same reduced matrix and the same residue character values.
Outside its stipulated integral half the indicator vanishes. The
frame and its integral inverse identify the integral half with the
support of the corresponding lattice indicator (the primes of the
denominator are excluded by the level hypothesis). Fourier inversion
thus expresses the test as a sum of vacuum translates with weights
$\widehat\phi(q)/|\mathcal Q|$. To check the sign, at a rational
input $u$ the character of the translate by the idempotent half of
$\delta\iota(c_{\rm sep}\widetilde q)/M$ is the inverse-transform
character of the level pairing. Indeed the pairing of the input
with $c_{\rm sep}\widetilde q$ is minus $\delta$ times the trace
of $\xi_0$ times the Hermitian pairing taken in the opposite order.
Alternation cancels this minus sign in the Fourier character.
Summing the characters over $q$ then gives $|\mathcal Q|$ on the
matching residue and zero on the others.

Applying $\omega_t(h,k)$ transports each translation by its
symplectic action. The finite representative of $\delta Y(q)$
differs from this moved translation by twice a vector of the moved
integral vacuum lattice for every $q$ with $\widehat\phi(q)\ne0$.
Here is the congruence calculation. Let $e$ be the tuple's finite
idempotent, $X$ the move by $k$ of the full rational translate, $u$
the finite representative of $\delta Y(q)$, $L_1$ the lattice after
the target move, and $L'$ the lattice after both moves. The split
support of $h$ is disjoint from the selected residue primes, so it
fixes the $e$ and $\bar e$ components. The marking gives
$u-X\in L_1$ by multiplying the rational columns by $M$ and using
the common residues. Fourier support gives $(1-e)X\in L_1$: the
character of a nonzero Fourier coefficient is trivial on every
period of the given residue function. Also $MX\in L_1$. Consequently
$u-eX\in L_1$. Put $v'=(u-eX)/2$. The class lattice pins give
$N L_1\subset L'$ and $k_0L_1\subset L'$ for an auxiliary $k_0$
prime to $M$, while $M\mid N$. Thus $2k_0 v'\in L'$. The primary
root condition gives $M^j u=2Nw$ with $w\in L_1$ for some $j\ge1$.
It follows that $M^jX\in L_1$ and
$M^j v'=Nw-e(e/2)M^jX\in L'$, since the selected primes
exclude $2$ and $eL_1\subset L'$. As $M$ is odd,
$\gcd(2k_0,M^j)=1$, so B\'ezout gives $v'\in L'$.

For the central character, choose $s$ in the indicated power of the
half-level ideal which is a unit on the conjugate idempotent, so
$\bar e/s$ is integral and $su=2Nw'$ with $w'\in L_1$.
Such $s$ exists by the coprimality of the two halves. Since
$e\bar e=0$, one has
$\bar e v'=N\bar e(\bar e/s)w'$. In the symplectic pairing with
$eX$ only this conjugate component contributes. Transfer the factor
$M$ to the first entry: $eMX$ and
$(N/M)\bar e(\bar e/s)w'$ both lie in $L'$. The additive character
is therefore one by self-duality. As the half-pairing against
$2v'$ is the full pairing against $v'$, the Heisenberg product law
and the doubled-lattice vacuum equation identify the two
translations with no residual scalar.

For precision about the size and phase of the projection, first
consider a single split support. On the target tensor lattice take
representatives of the new lattice modulo its intersection with the
old. The finite sum of translates weighted by the new lattice
semicharacter, divided by their number $J$, is precisely the
new vacuum projection. Distinct representatives give orthogonal
translates of the unit old vacuum: self-duality separates their
difference by a character of the old lattice. Their unnormalized
sum consequently has squared norm $J$. In the split model the
Weil operator has the phase $\chi_W(y)^{-1}$ times a positive
multiple of that sum. This is the positive split intertwiner of
CW6: the two complementary halves give finite translate sums with
positive inverse square-root normalization, and averaging against
the new semicharacter leaves that same phase. Unitarity now gives
the positive multiple $J^{-1/2}$. Thus on the classical reading the
factor is $\chi_W(y)^{-1}\sqrt J$ times the normalized projection.
Both the old and new rational translates read their specified
classical Heisenberg translates: the finite representative of
the negative row vector of $z$ reads translation by $z$, with
the lattice semicharacter unchanged under negation.

For disjoint split supports the intermediate tensor lattice $B$
is the sum of $B\cap A$ and $B\cap C$, with $A,C$ the two
end lattices. This holds on each local factor since there one of
the two moves is the identity, and hence on the global lattices.
In the two successive finite averages the classes from $B\cap A$
act on the starting vacuum by its prescribed character and the
classes from $B\cap C$ are absorbed in the final average. All
resulting classes have the same number of preimages. Thus the
composition of normalized projections on the original vacuum
equals the final projection. The two relative-index factors
multiply, as is also seen by taking local indices at disjoint
supports. Induction gives the factor for the full support.
The relative index in the tensor lattice is $I^r$: locally the
target lattice is free of rank $r$, and indices multiply over
places. The determinants multiply, giving the one character
$\chi_W(y)^{-1}$ for $h$.

Pair this result with a homogeneous creation polynomial $P$ of degree
$d$. In Gaussian coordinates a creation operator on the reading is
the holomorphic direction operator with its Gaussian connection
term. At the origin the evaluation of the normalized creation
polynomial gives $(\sqrt{2\pi\xi_R})^{-d}$ times the classical
theta derivative. For a letter with direction $u$, differentiating the real
Heisenberg translation in directions $u$ and $iu$ gives on its
holomorphic factor the operator
\[
 (2\pi\xi_R)^{-1/2}\!
 \left[\tfrac12(D_u-iD_{iu})
            -\pi\sum_b H_x(u_b,z_b)\right].
\]
Here $D$ is the real directional derivative. Each $H_x(u_b,z_b)$
is antiholomorphic linear in $z_b$, so the holomorphic derivative
kills it, and its value at the origin is zero. Iterating on a
monomial and evaluating at the origin leaves only the classical
derivative, with one factor $(2\pi\xi_R)^{-1/2}$ per letter.
This is the same direction convention as the
classical harmonic reading. The finite translation of $\delta Y$
reads the classical translate by $-\delta Y$. Reflection fixes the
canonical even vacuum and commutes with the lattice projection.
Replacing all roots by their negatives replaces $z$ by $-z$ in the
classical theta function. A homogeneous order-$d$ derivative at zero
therefore acquires $(-1)^d$, so the effect of $-\delta$ is
$(-1)^d\delta^d$. By linearity the resulting value is
\[
 \frac{\delta^d(-1)^d}{(\sqrt{2\pi\xi_R})^d}\,
 b\,\chi_W(y)^{-1}\,\frac{\sqrt{I^r}}{|\mathcal Q|}
 \ \mathrm{thP}\!\left(
    \mathrm{gphi}(\Lambda,L,v,\widehat\phi,Y),P,x_k\right).
\]
Here $\mathrm{gphi}$ forms the indicated sum of classical
Heisenberg translates of the projected vacuum, and $x_k$ is the
moved frame. This is the one-class formula $\mathrm{ClsAt}$ on
the original lattices and characters.
% END P74I6

% BEGIN P74I7
\paragraph{P74I7. One point function and the identity level element.}
We now prove \texttt{point\_identity\_core}. If there is no pinned
record use the constant function $1$ and $z=1$. Otherwise choose
one pinned record for P74I5. Its lower model, reference frame and
dictionary come from $T$ at the distinguished embedding. In any
other pinned record $O$, restriction of its coefficient embedding to
the lower field is that same distinguished embedding. Its lower
dictionary and $\mu$ are transported along the given identification
$T.V=W$. Thus for each adelic $g$ its component unit agrees with
$C_{k(g)}$ of the anchor, where $k(g)$ is the component in the
anchor dictionary. This transport consists in substituting equal
embeddings and equal Hermitian spaces. The tables of higher
components are allowed to differ. Apply P74I3 for the fixed lower
coefficient map, and P74I6 for the tied lattice isometry. Put
\[
 d=\sum_{\tau}\sum_{j<\mathrm{mult}(\tau)}\mathrm{wtA}(\kappa_i)(\tau,j),
 \qquad K_d=a\delta^d(-1)^d(\sqrt{2\pi\xi_R})^{-d}.
\]
On $\mathrm{AutQ}(W)$, using the representative $q.\mathrm{out}$
specified by the point reader, set
\[
           \lambda(q)=K_d C_{k(q.\mathrm{out})},\qquad z=1.
\]
These are independent of $O,t,c$. The factors $a$ and $C_k$ have
modulus one, $\delta$ is an integer unit and $\xi_R>0$ is the
absolute value of the distinguished embedding of the nonzero
imaginary parameter. Hence $\lambda$ is everywhere nonzero.

Take $t$, placed coefficients $x$, an original section $f$ with
$\mathrm{IsCls}(O,t,x,f)$, and $g=q.\mathrm{out}$. The placement
pin supplies a definite form $y$ whose values are exactly the
coefficient readings $O.c'_0(x_j)$ on the prescribed representatives
$h_j$, and $O.\mathrm{thAd}$ is the reading of its theta lift.
Use the rational translate $\gamma g$ in the dictionary, a common
compatible pair of lower and higher markings, and its determinant
root $x_\gamma$. Self-duality of the period lattice gives the
dual roots $Y_j$ used by both short-series formulas. Write
$\kappa'$ for the factor taking the classical vacuum theta to the
product of the lower generator readings in this chart. Let $S(g)$
be the finite class sum in this common marking and frame, with
these $Y_j$, the stored Fourier coefficients, and $O.\mathrm{io}(x_j)$.

On the geometric side the marked lift and the section's class
value, as in T73C1--T73C3, evaluate the actual $f$ as
\[
       \mathrm{rdU}(f;O.\mathrm{dict}(g))
        =c\,\xi_t(x_\gamma)\,\kappa' S(g)
\]
by P74I4. On the adelic side unfold the compact definite quotient
into its specified double cosets. The weight of coset $j$ is the
given subgroup volume divided by its stabilizer size $\gamma_j$.
Indeed its quotient measure is that subgroup volume modulo the
finite stabilizer, and the opposite level characters of kernel and
form cancel on the subgroup. P74I3 supplies the factor $a$. P74I5
reads the base moved vacuum as
$C_{k(g)}O.\mu(g)\kappa' v$. The twist from $t_*$ to $t$ multiplies
this by $\xi_t(x_\gamma)$. This follows by twisting the Weil action
on the determinant root and using its uniqueness on theta functions. The base character is trivial. Apply P74I6 to each coset.

For clarity, if $m_j$ and $Z_j$ denote the lattice's stored
$\mathrm{mIdx}$ and $\mathrm{zCard}$ at $N_j$, its arithmetic
factor satisfies
\[
  \iota_O(\mathrm{ep}_j Z_j/m_j)
     =\chi_W(y_j)^{-1}\sqrt{[\Lambda_j:\Lambda_j\cap S]^r}.
\]
Here $S=\mathcal O_E^n$ and $y_j$ is the supported determinant
of $h_j$. The ideal of its supported finite idele gives the
class lattice's determinant up to conjugation. Outside its support
the characters are unramified or trivial on units, by the bad-prime
pin. The original tuple ideal-character pin gives its character
value, and the squared norm pin for this ideal supplies the positive
square root of the index (RTG9). Conjugating the class lattice
preserves its index against $S$. Substitution and the definition of
$\mathrm{ep}_j$ give the displayed equality, with the same
$Z_j,m_j$ as in the geometric class sum. The Fourier coefficient
of the finite test is the embedded stored weight by the tuple
Fourier pin (RTG11). Thus summing the formula of P74I6 gives
\[
 \mathrm{thF}_t(y)(g)=K_d C_{k(g)}O.\mu(g)\,
                         \xi_t(x_\gamma)\kappa' S(g).
\]
Multiplying by $c$ and substituting the geometric value proves the
identity even if $S(g)$ or $c$ vanishes. For $z=1$ the full
level-action laws give $\mathrm{actLev}(1)f=f$, so the product
$O.\mu(g)\mathrm{rdU}(f;O.\mathrm{dict}(g))$ is exactly
$\mathrm{rdQ}(O,1,f,q)$. For every tuple choose the root of unity
$1$ in the distinguished original field, of order one. We have
proved for all the data in the original quantifier order that
\[
      c\,O.\mathrm{thAd}(t,x,q)
        =\lambda(q)\,\mathrm{rdQ}(O,1,f,q).
\]
The section and the coefficient lift here are the ones given by the
test. Neither surjectivity of a space of sections nor spanning by
class sums is needed for this identity.
% END P74I7
